\documentclass[11pt]{article}
\usepackage[T1]{fontenc}
\usepackage{lmodern}
\usepackage[margin=1in]{geometry}
\usepackage{amsmath,amssymb,amsthm,microtype}
\usepackage[colorlinks=true,linkcolor=blue,citecolor=blue,urlcolor=blue]{hyperref}
\newtheorem{theorem}{Theorem}
\newtheorem{lemma}[theorem]{Lemma}
\newcommand{\kk}{\Bbbk}
\newcommand{\supp}{\operatorname{supp}}
\newcommand{\LL}{\mathcal L}
\title{Minimal persistence grids and homological Morse bounds}
\author{Kaiwen Jin\\\small\href{mailto:yui.ui.siki@gmail.com}{yui.ui.siki@gmail.com}}
\date{October 14, 2026}
\begin{document}
\maketitle

\begin{abstract}
The persistent homology of a finite simplicial complex is determined by a
canonical grid whose coordinate values occur in minimal generators and
relations. We prove that the number of values in each coordinate is bounded
by the homological Morse lower bound defined from the multigraded Betti
tables. The result holds over every field, for every number of parameters,
and for filters with coincident coordinate values. It establishes the
conjecture in Remark~5.12 of Harrington, Tillmann and Torras-P\'erez.
The proof selects a minimal presentation degree on each coordinate slice.
At that degree a local Koszul calculation eliminates higher syzygies;
injectivity of top-dimensional homology handles the final homology degree.
\end{abstract}

\section{The grid and the lower bound}

Persistent homology records the homology of sublevel complexes together
with the maps induced by their inclusions. In the inverse problem, one asks
which filters give the same persistent homology. Harrington, Tillmann and
Torras-P\'erez~\cite{HHT} obtain dimension estimates for these spaces of
filters using both a minimal determining grid and a lower bound for
homological Morse numbers. Their Remark~5.12 asks whether each coordinate
of the minimal grid has cardinality at most that lower bound. We prove
this comparison for general filters, including coordinate coincidences.

Fix a field $\kk$, an integer $n\ge1$, and a finite nonempty simplicial
complex $K$ of dimension $d$. A filter assigns $f(\sigma)\in[0,1]^n$
to every nonempty simplex, with $f(\tau)\le f(\sigma)$ coordinatewise
whenever $\tau\subseteq\sigma$. For $u\in\mathbb R^n$, ordered
coordinatewise, set
\[
 K_u=\{\sigma\in K:f(\sigma)\le u\},\qquad
 M_u^q=H_q(K_u;\kk),\qquad 0\le q\le d,
\]
using ordinary unreduced homology and inclusion-induced maps. Write
$\mathbf M=(M^0,\ldots,M^d)$. A free summand born at a grade $w$
has value $\kk$ at every $v\ge w$ and zero elsewhere; all maps
between its nonzero components are identities.
Let $\xi_i^q(u)$ be the multiplicity of a free summand at grade $u$
in position $i$ of a minimal free resolution of $M^q$. We set these
numbers to zero when $q<0$ or $q>d$.

Let
\begin{equation}\label{eq:grid}
 S=\bigcup_{q=0}^d\supp(\xi_0^q+\xi_1^q),\qquad
 G_j=\{s_j:s\in S\},\qquad G=\prod_{j=1}^nG_j.
\end{equation}
By \cite[Lemma~1.23 and Definition~2.9]{HHT}, $G$ is the canonical
minimal grid determining the tuple. Values $0$ and $1$ are included only
when they occur in the indicated supports.
The lower bound in question is
\begin{align}
 b_q(u)&=\max\left\{0,\xi_0^q(u)+\xi_1^{q-1}(u)
       -\sum_{p=1}^{n-1}\xi_{p+1}^{q-p}(u)\right\},\label{eq:b}\\
 \LL(\mathbf M)&=\sum_{u\in G}\sum_{q=0}^d b_q(u).\label{eq:L}
\end{align}
The expression in \eqref{eq:b} comes from the Morse inequalities of
Guidolin and Landi~\cite[Theorem~7.3]{GL}; \eqref{eq:b}--\eqref{eq:L}
are the quantities used in \cite[Definition~5.7 and Proposition~5.9]{HHT}.
Their grid comparison is known for one parameter and for filters
injective in every coordinate~\cite[Lemma~5.10 and Remark~5.12]{HHT}.
For general filters the higher-syzygy terms can cancel the positive
terms at some grades. We will find a grade with a positive contribution
on every coordinate slice of $G$.

We use consecutive integer indices for each ordered $G_j$ and extend
the modules along this grid to $\mathbb Z^n$. This regrading preserves
minimal free resolutions and their summand multiplicities
\cite[Section~1.4]{HHT}. Each coordinate set keeps its cardinality. In these indices put
$R=\kk[x_1,\ldots,x_n]$, with $\deg x_j=e_j$, where $e_j$ is the
$j$th standard basis vector. A free summand $R(-w)$ has a copy of $\kk$
at $v$ precisely when $w\le v$, and its maps between present components
are identities. Each module has a minimal presentation
\begin{equation}\label{eq:presentation}
 F_1^q\longrightarrow F_0^q\longrightarrow M^q\longrightarrow0,
\end{equation}
whose basis degrees are counted by $\xi_1^q$ and $\xi_0^q$.
The common grid determines every component, even if a component has a
smaller minimal grid. Below its support the extended module is zero.

\section{A local Koszul calculation}

For a finitely presented integer-graded $R$-module $V$ and a grade $u$,
its Koszul complex is
\[
 \mathbb K_i(V)(u)=
 \bigoplus_{\substack{\alpha\subseteq\{1,\ldots,n\}\\|\alpha|=i}}
 V_{u-e_\alpha},\qquad e_\alpha=\sum_{k\in\alpha}e_k.
\]
The differential is the alternating sum of the coordinate maps that
delete one index from $\alpha$. Its homology computes the Betti tables:
\begin{equation}\label{eq:Koszul}
 \xi_i(u)=\dim_\kk H_i(\mathbb K_*(V)(u));
\end{equation}
see the integer-graded derivation in \cite[Section~1.4]{HHT}, followed
by Definition~1.30 and Lemma~1.31 there. The next lemma needs
isomorphisms in just one coordinate, away from the vertex $u$.

\begin{lemma}\label{lem:local}
Fix $j\in\{1,\ldots,n\}$. Suppose that
\begin{equation}\label{eq:off-center}
 V_{u-e_j-e_\alpha}\longrightarrow V_{u-e_\alpha}
\end{equation}
is an isomorphism for every nonempty
$\alpha\subseteq\{1,\ldots,n\}\setminus\{j\}$.
Then $\xi_i(u)=0$ for all $i\ge2$. If, in addition,
$V_{u-e_j}\to V_u$ is injective, then $\xi_1(u)=0$.
\end{lemma}

\begin{proof}
Let $A$ and $B$ be the Koszul complexes in the variables other than
$x_j$, at grades $u-e_j$ and $u$, respectively. Multiplication by $x_j$
gives a chain map $\phi:A\to B$. The hypothesis says that
$\phi_p$ is an isomorphism for every $p\ge1$, including zero terms
above the lengths of these complexes.

Choose the exterior ordering with $j$ first. Separating the subsets
that contain $j$ identifies the full Koszul complex with
\[
 C_i=B_i\oplus A_{i-1},\qquad
 D_i(b,a)=(d_Bb+\phi_{i-1}a,-d_Aa).
\]
The minus sign comes from differentiating an exterior product
$e_j\wedge\eta$: its two terms are $x_j\eta-e_j\wedge d\eta$.

For a cycle $(b,a)\in C_i$ with $i\ge2$, take
$c=\phi_i^{-1}b$. The cycle equation and $d_B\phi=\phi d_A$ give
\[
 \phi_{i-1}(a+d_Ac)=d_Bb+\phi_{i-1}a=0.
\]
Since $\phi_{i-1}$ is injective, $a=-d_Ac$. Thus
\[
 (b,a)=D_{i+1}(0,c)
\]
is a boundary, proving the first assertion by \eqref{eq:Koszul}.

For a degree-one cycle, the same choice
$c=\phi_1^{-1}b$ reduces it, modulo the boundary $D_2(0,c)$, to
$(0,a+d_Ac)$ with $\phi_0(a+d_Ac)=0$. Injectivity of $\phi_0$
makes that cycle zero. When $n=1$, $A$ and $B$ have only degree-zero
terms, so the assertion is the kernel calculation for
$V_{u-e_j}\to V_u$.
\end{proof}

\section{The cardinality bound}

The last homology degree requires a simple fact about the common
underlying complex. For every inclusion $K_v\subseteq K_w$, the map
$H_d(K_v;\kk)\to H_d(K_w;\kk)$ is injective: there are no
$(d+1)$-simplices, so these homology groups are the kernels of the
degree-$d$ boundary, and the chain map includes one kernel in the
other. At $d=0$, ordinary unreduced homology has $\partial_0=0$,
and the same argument applies. Thus all coordinate maps of $M^d$
are injective, also after the grid regrading.

\begin{theorem}\label{thm:main}
Let $K$ be any finite nonempty simplicial complex, $\kk$ any field,
and $n\ge1$. For every face-monotone filter $f:K\to[0,1]^n$,
including filters with coincident coordinate values, its ordinary
unreduced persistent homology tuple satisfies
\[
 |G_j|\le
 \sum_{u\in G}\sum_{q=0}^{d}
 \max\left\{0,\xi_0^q(u)+\xi_1^{q-1}(u)
             -\sum_{p=1}^{n-1}\xi_{p+1}^{q-p}(u)\right\}
 =\LL(\mathbf M)
 \quad (1\le j\le n).
\]
Here $d=\dim K$, $\xi_i^q$ are the multigraded Betti tables of
$H_q(K_u;\kk)$ from this same filter, with absent homology degrees
assigned zero, and $G=\prod_jG_j$ is the canonical minimal grid
defined in \eqref{eq:grid}.
\end{theorem}

\begin{proof}
Fix $j$ and $t\in G_j$. In the finite nonempty set
$\{s\in S:s_j=t\}$ choose an element $u$ minimal in the
coordinatewise order. There is consequently no $s\in S$ with
$s_j=t$, $s\le u$ and $s\ne u$. This choice uses the
presentation degrees from all homology degrees together.

For a nonempty
$\alpha\subseteq\{1,\ldots,n\}\setminus\{j\}$ put
$v=u-e_\alpha$. No basis degree $w$ of either $F_0^q$ or $F_1^q$
can satisfy $w\le v$ and $w_j=t$, for such a $w$ would belong to
$S$ and be strictly below $u$ on the same slice. All free summands
present at $v$ therefore have $w_j\le t-1$, and are already present
at $v-e_j$. Multiplication by $x_j$ gives isomorphisms on both
free modules of \eqref{eq:presentation} at these two grades.
Their commuting presentation maps identify the images, hence the
cokernels. We obtain
\[
 M^q_{u-e_j-e_\alpha}\ \xrightarrow{\ \cong\ }\ M^q_{u-e_\alpha}
 \qquad(0\le q\le d).
\]
Lemma~\ref{lem:local} yields
\[
 \xi_i^q(u)=0\quad(i\ge2,\ 0\le q\le d).
\]
Its last assertion, applied to the injective coordinate map of
$M^d$, also gives $\xi_1^d(u)=0$.

Since $u\in S$, some $\xi_0^q(u)$ or $\xi_1^q(u)$ is positive.
In the latter case $q<d$. At this selected grade all subtraction
terms in \eqref{eq:b} vanish, so
\[
 b_q(u)=\xi_0^q(u)+\xi_1^{q-1}(u),\qquad
 \sum_{q=0}^d b_q(u)\ge1.
\]
Each selected $u=u(t)$ belongs to $G$. Distinct values of $t$ give
distinct grades because $u(t)_j=t$. Nonnegativity of every $b_q$
now gives
\[
 \LL(\mathbf M)
 \ge\sum_{t\in G_j}\sum_{q=0}^d b_q(u(t))
 \ge |G_j|.
\]
This proves the result for each coordinate. An empty coordinate
set would give the inequality $0\le\LL(\mathbf M)$ directly.
\end{proof}

\paragraph{AI-use statement.}
OpenAI Codex assisted with proof development, source checking and
manuscript preparation. Independent mathematical review is requested.

\begin{thebibliography}{9}
\bibitem{HHT}
H.~A.~Harrington, U.~Tillmann and M.~Torras-P\'erez,
\emph{The fiber of multiparameter persistent homology for simplicial
complexes}, preprint, 2026, arXiv:2608.04640v1.
\url{https://arxiv.org/abs/2608.04640v1}.

\bibitem{GL}
A.~Guidolin and C.~Landi,
\emph{Morse inequalities for the Koszul complex of multi-persistence},
Journal of Pure and Applied Algebra \textbf{227} (2023), no.~7, 107319.
\href{https://doi.org/10.1016/j.jpaa.2023.107319}{doi:10.1016/j.jpaa.2023.107319}.
\end{thebibliography}
\end{document}
